\documentclass[a4paper,12pt]{article}

% Pacotes básicos
\usepackage[utf8]{inputenc}
\usepackage[T1]{fontenc}
\usepackage[brazil]{babel}
\usepackage{amsmath, amssymb}
\usepackage{geometry}
\geometry{margin=2.5cm}

% Pacote para códigos em Python
\usepackage{listings}
\usepackage{xcolor}
\renewcommand{\lstlistingname}{Algoritmo}
% Configuração do ambiente listings para Python
\lstset{
    language=Python,
    basicstyle=\ttfamily\small,
    keywordstyle=\color{blue}\bfseries,
    stringstyle=\color{red},
    commentstyle=\color{gray}\itshape,
    numbers=left,
    numberstyle=\tiny\color{gray},
    stepnumber=1,
    numbersep=5pt,
    backgroundcolor=\color{white},
    frame=single,
    tabsize=4,
    showspaces=false,
    showstringspaces=false,
    breaklines=true,
    breakatwhitespace=true,
    captionpos=b
}

\begin{document}

    \begin{enumerate}
        \item Dada uma função linear $f:\mathbb{R}^n \to \mathbb{R}^m$ e a dimensão do vetor de entrada $n$, escreve um pseudo código que retorne a matriz $A$ tal que $f(x) = Ax$.\\
        \textit{Dica: caso queira abstrair um método que crie um vetor ou matriz específica, apenas explique o que esse método faz. Exemplo: \texttt{cria\_vetor\_zero(n)} cria um vetor de zeros de dimensão $n$.}

        \item Sobre o seguinte algoritmo:
        \begin{lstlisting}[gobble=12,caption=Exemplo de código Python]
            novos_vetores = []
            for v_k in vetores:
                soma = 0
                for u_i in novos_vetores:
                    soma += u_i.T @ v_k / (u_i.T @ u_i) * u_i
                u_k = v_k - soma
                norm = np.linalg.norm(u_k)
                if np.abs(norm) < 1e-10:
                    print("conjunto LD")
                    continue
                u_k = u_k / (norm)
                novos_vetores.append(u_k)
        \end{lstlisting}
            \begin{enumerate}
                \item Explique a linha 5 do código.
                \item Qual a finalidade do algoritmo?
            \end{enumerate}
            \item A temperatura T de um dispositivo com três processadores é afim da potência \(\mathbf{P}=(P_1,P_2,P_3)\). Dados:
            $$
            T(10,10,10)=30,\quad T(100,10,10)=60,\quad T(10,100,10)=70,\quad T(10,10,100)=65.
            $$
            Explique como o código a seguir pode ser usado para determinar a função $T(P) = a^T P + b$, onde $a \in \mathbb{R}^3$ e $b \in \mathbb{R}$.
            \begin{lstlisting}[gobble=16]
                A = np.array([
                    [10, 10, 10, 1],
                    [100, 10, 10, 1],
                    [10, 100, 10, 1],
                    [10, 10, 100, 1]
                ])
                b = np.array([30, 60, 70, 65])
                x = np.linalg.solve(A, b)
                a = x[:3]
                b = x[3]
            \end{lstlisting}
    \end{enumerate}

\end{document}